\section{Discussion}
\label{sec:discussion}

\subsection{The Mechanistic Claim and the Null Result: Two Contributions}

The h-e1 and h-m1 results are complementary contributions at different levels of the causal chain.

\textbf{h-e1 establishes the mechanistic claim} at the highest confidence level: a mathematical guarantee. Ratio reward produces non-zero advantage variance in any GRPO group where at least one completion achieves partial credit and no completion achieves full credit. This guarantee is independent of model scale, dataset, and training configuration.

\textbf{h-m1 establishes the prerequisite} for observing the mechanistic guarantee in practice. The null result is a principled finding: it identifies the binding constraint (generation length must allow syntactically complete solutions) and rules out the hypothesis that ratio reward helps regardless of experimental setup. Together, h-e1 and h-m1 answer different questions with equally high confidence: ``Does the mechanism exist?'' (yes, mathematically) and ``Was the mechanism active in our policy-level experiment?'' (no, and here is exactly why).

\subsection{Why the Null Result Is Not a Failed Experiment}

The h-m1 null result is not ambiguous. The FractionPartialCallback diagnostic provides a real-time monitor of the partial-pass prerequisite, and \texttt{clipped\_ratio = 1.0} at every step provides a direct measurement of the root cause. We know exactly why both rewards are equivalent. This precision has practical value: our diagnostic framework (FractionPartialCallback + clipped\_ratio monitoring) provides a template for detecting this failure mode in future RLEF experiments, preventing wasted compute on experiments where the reward function choice is irrelevant.

\subsection{Honest Limitations}

\textbf{Policy-level claims are unverified.} The original hypothesis (ratio reward improves HumanEval pass@1 by $\geq$3pp) could not be evaluated. We make no claim about ratio reward's benefit for HumanEval, APPS all-pass rate, or cross-benchmark generalization. These remain open empirical questions requiring correctly-powered experiments.

\textbf{Single model, single dataset.} All experiments use DeepSeek-Coder-6.7B-instruct on APPS. The mechanistic claim (h-e1) is model-agnostic, but the policy-level analysis (h-m1) covers only one model-dataset combination.

\textbf{Gradient norm analysis covers steps 1--136 of h-m1.} The gradient norm 95\% bootstrap CI on (ratio $-$ binary) over steps 1--136 includes zero (mean = $+$0.000730, CI = $[-$0.000253, $+$0.002561]), consistent with both conditions receiving identical zero-reward signals. Longer training or correctly-powered conditions may exhibit different patterns.

\subsection{Implications for RLEF Practitioners}

Pre-training checklist for ratio reward experiments: (1) validate base solve rate --- compute \texttt{reward\_mean} and \texttt{fraction\_partial} for a small sample under the planned generation length; if both = 0, ratio and binary rewards are equivalent; (2) monitor \texttt{clipped\_ratio} during training --- if $\to$ 1.0, no reward signal reaches the model regardless of reward formulation; (3) use ratio reward when \texttt{fraction\_partial} $> 0$ at early steps; (4) for hard datasets, increase \texttt{max\_new\_tokens} to 1,024--2,048 or switch to easier training data (HumanEval, MBPP, APPS introductory split).

\subsection{Broader Impact}

Our analysis quantifies a previously unexamined failure mode in GRPO-based RLEF. The one-line fix (switch from binary to ratio reward) is immediately actionable. The diagnostic framework adds minimal overhead and provides real-time visibility. We release all code, configurations, and diagnostic tools to support reproducibility.
